\section{Introduction}

Di Paola and Heller introduced dominical categories to develop recursion
theory in categorical algebra. They sought a formulation of generalized
incompleteness in this setting \cite[pp.~594--595]{DPH}.
We construct internal deductive systems in completions of recursion
categories and apply categorical inseparability arguments to their
sentence theories. The resulting essential-incompleteness theorems
assert that every consistent extension in a specified effective class
is incomplete. The three principal applications concern arithmetic,
a theory of real closed fields expanded by constants, and propositional
theories obtained from cosheaves.

Di Paola and Heller observed that diagonalizable algebras and the modal
logic $\mathrm{GL}$ include the formalized L\"ob theorem among their special
axioms, so that second incompleteness is largely built into these systems.
They related this axiom to arithmetic provability formulas satisfying
the Hilbert--Bernays derivability conditions. They also discussed
difficulties in extending modal fixed-point theorems to numerals and
quantification. Their objective called for an intrinsic recursion theory
supplying the fixed-point constructions of recursion theorems and
analogues of creative and effectively inseparable sets
\cite[pp.~595--596]{DPH}.

The connection between deduction and recursion uses two representations
of the same sentence theory: a subobject of a sentence object and a
restriction idempotent on that object. For partial functions, a restriction
idempotent is a partial identity on a subset of the source; splitting it
makes that subset an object. Let $\mathscr R$ be a recursion category.
Its \emph{restriction idempotent completion}, written
$\mathscr S=\Split(\mathscr R)$, is obtained by splitting the restriction
idempotents of $\mathscr R$. The \emph{category of total morphisms} of
$\mathscr S$, written $\mathcal D=\Tot(\mathscr S)$, has all objects
of $\mathscr S$ and only its total morphisms. A morphism is total when its restriction
idempotent is the identity on its source.
In the applications, $\mathcal D$ is regular and supports the syntax
and initial algebras required in Appendix~\ref{app:syntax-theories}.
For first-order logic this includes terms, formulas, quantifiers,
equality, and admissible substitution. The application sections prove
that their realizations have the stated universal properties.

Let $\Phi\in\mathcal D$ be the formula object, let
$j:S\hookrightarrow\Phi$ be its sentence subobject, and choose
an axiom subobject $i_A:A\hookrightarrow S$.
The inference rules and axioms determine an initial algebra of finite
derivations with conclusion morphism $c_A:P_A\to\Phi$.
Pulling back along $j$ gives the sentence-conclusion morphism
$c_{A,S}:P_{A,S}\to S$. Its regular image
\[
 P_{A,S}\twoheadrightarrow T_{A,S}\xhookrightarrow{m_A}S
\]
is the sentence part of the deductive closure of the axioms
(Theorem~\ref{thm:first-order-theory}). It represents the sentences
derivable from $A$ under the internal inference rules.

When $m_A$ admits a partial inverse in $\mathscr S$, this subobject
also represents the restriction range
$t_A=\ran(c_{A,S}):S\rightharpoonup S$
(Lemma~\ref{lem:range-regular-image}).
For the total sentence-negation morphism $\nu:S\to S$,
the inverse image $\nu^*t_A$ is the restriction idempotent representing
refutable sentences.
Consistency means $t_A\nu^*t_A=0_S$, and completeness means
$t_A\vee\nu^*t_A=1_S$. These equations put the logical
questions in terms of restriction idempotents.

There are two incompleteness arguments. The first uses a range functor
$F:\mathscr R\to\mathcal A$ into a split restriction category,
extended to $\widehat F:\mathscr S\to\mathcal A$ by chosen
splittings. When $\mathcal A$ is the partial-map category of a regular
category, a sound interpretation factors the interpreted conclusion morphism
through a validity subobject. If validity is disjoint from
its inverse image under negation, the interpreted restriction idempotents
representing theorems and refutations are disjoint (Proposition~\ref{prop:interpreted-theory}).
A total reduction from the pair of restriction idempotents determined by
a source Turing morphism, together with a source lift of a specified partial conditional in the
target, then yields incompleteness of the source theory
(Theorem~\ref{thm:logical-comparison}).
Without the lifting hypothesis, this pair can be
separable even when the range functor is faithful, preserves restriction
products, and reflects zero morphisms
(Appendix~\ref{app:lifting-counterexample}).

The second argument starts with an inseparable pair in the source.
A total morphism reduces that pair to the restriction idempotents
representing theorems and refutations of a base theory. Enlarging the axiom
subobject enlarges the generated theory, so the same reduction applies to each represented extension.
If the extension is consistent, a separator of the corresponding restriction
idempotents would pull back to a separator of the fixed pair.
Corollary~\ref{cor:domain-incompleteness} therefore gives inseparability,
and hence incompleteness, for every such consistent extension.
The realizations specify which effective axiom sets are represented.

The first-order syntactic realization uses the recursion categories
$S_T$ of Montagna and Di Paola--Montagna \cite{Montagna,DPM}, for
consistent computably enumerable (c.e.) extensions $T$ of Peano arithmetic
($\PA$). Their morphisms are equivalence classes of indices of partial
recursive functions under universally $T$-provable graph equality.
We realize the syntax and derivation algebras in
$\mathcal B_T=\Tot(\Split(S_T))$, with their universal properties
expressed by this equality
(Theorem~\ref{thm:syntactic-generated-theory}).
For the arithmetic application we fix the ambient category $S_{\PA}$
and vary the represented axiom subobjects; the represented theory need
not be the theory that defines equality of morphisms.

Soundness in this application separates the interpretation of syntax
codes from the interpretation of the formulas they encode.
Standard arithmetic interprets the codes, while a model $K$ of the
represented axioms interprets the decoded formulas.
Induction on finite derivations gives the soundness factorization
(Proposition~\ref{prop:standard-code-soundness});
Theorem~\ref{thm:ultrapower-soundness} transports it to arithmetic
ultrapowers. Partial proof search supplies the lift needed for the
comparison theorem. We obtain incompleteness in the fixed syntactic
category and in its standard and ultrapower interpretations
(Theorem~\ref{thm:arithmetic-internal-comparison}).
Every consistent c.e.\ extension of PA has recursively inseparable
theorem and refutation sets (Corollary~\ref{cor:arithmetic}).
We also construct an independent sentence uniformly from an index for
an axiom enumeration of any such theory
(Theorem~\ref{thm:uniform-independent-sentences}).
The numerical realization supplies a second arithmetic argument and
is used in the cosheaf application.

The Miller application uses a sequence $a:\N\to[0,1]$ constructed
by Miller \cite{Miller2004}. A morphism of the recursion category
$\mathcal E_a$ is a partial numerical function computed by one oracle
program with the same domain and values for every binary representation
of $a$. First-order logic with equality is realized in
$\mathcal D_a=\Tot(\Split(\mathcal E_a))$. The partial identities on
\[
 L_a=\{n:a(n)<1/2\},\qquad R_a=\{n:a(n)>1/2\}
\]
are inseparable in $\mathcal E_a$, although each fixed binary
representation computes a separator.
The theory $T_a$ of real closed fields expanded by constants $c_n$
contains the strict rational bounds satisfied by $a(n)$.
It is consistent, and the total morphism sending $n$ to the sentence
$c_n<1/2$ reduces this inseparable pair to the restriction idempotents
representing theorems and refutations. Every consistent extension is
incomplete if a single oracle program semidecides its axiom set using
any binary representation of $a$ as oracle
(Theorem~\ref{thm:miller-incompleteness}).

The third application constructs a recursion category $\mathscr R_X$
from cosheaves on a three-point space $X$.
It is not equivalent as a restriction category to the category of partial
recursive functions. Its completion realizes propositional syntax and
finite derivations by initial algebras. The generated theory is a
subfunctor of the sentence functor; evaluation at each of the two endpoints
gives a consistent propositional theory.
Both theories are essentially incomplete: every consistent extension
in the same language with a c.e.\ axiom set is incomplete
(Corollary~\ref{cor:cosheaf-essential-incompleteness}).

Section~\ref{sec:lifting} proves the comparison theorem and the criterion
using source conditionals, after the categorical preliminaries.
Section~\ref{sec:internal-theories} applies these results to the
sentence theories constructed in Appendix~\ref{app:syntax-theories}.
Sections~\ref{sec:syntactic} and~\ref{sec:models} establish the syntactic
realization and soundness; Sections~\ref{sec:arithmetic},
\ref{sec:miller}, and~\ref{sec:cosheaves} give the three applications.
The remaining appendices supply the cosheaf infrastructure, the lifting
counterexample, and the auxiliary faithful model family and its behavior
under theory extension.

\subsection{Related work}

Di Paola and Heller proposed inseparability of theorems and refutable
sentences as part of a categorical formulation of incompleteness
\cite[pp.~594--595]{DPH}. Local connectedness enters Heller's existence
theorem for recursion categories \cite[Sections 8--9]{Heller}.
Di Paola and Montagna's syntactic examples led to the question of
extending that theorem beyond locally connected categories \cite{DPM}.
Lengyel proved existence theorems using faithful iteration functors
\cite{Lengyel2004,Corrections}.
The comparison for internal theories here starts with a range functor
$F:\mathscr R\to\mathcal A$ from a recursion category $\mathscr R$
to a split restriction category $\mathcal A$. These are the source and
target categories of $F$, respectively; the Turing morphism in
$\mathscr R$ is part of the hypotheses.

The extension $\widehat F:\Split(\mathscr R)\to\mathcal A$
interprets the internal sentence theories in the target category.
When the validity subobject is disjoint from its inverse image under negation,
soundness gives disjoint interpreted theorem and refutation restriction
idempotents (Proposition~\ref{prop:interpreted-theory}). Under the reduction and lifting
hypotheses of Theorem~\ref{thm:logical-comparison}, this gives
incompleteness both of the interpreted theory in $\mathcal A$ and
of the internal theory in $\Split(\mathscr R)$.
For the arithmetic comparison, let $M$ be the standard model of
arithmetic on $\N$ or an ultrapower of that model. Interpreting program
graphs in $M$ defines the range functor $F_M:S_{\PA}\to\Pfn$,
where $\Pfn$ is the category of sets and partial functions.
Appendix~\ref{sec:family} fixes a consistent theory $T\supseteq\PA$
and the family $\mathcal M_T$ of its arithmetic models with underlying
set $\N$. The range functor
$F_T:S_T\to\prod_{M\in\mathcal M_T}\Pfn$ combines their
program-graph interpretations.

Rosolini's $P$-categories generalized dominical categories
\cite{RosoliniThesis,Rosolini}. Heller adopted them and defined
prodominical categories \cite[Section 5]{Heller}.
Cockett and Lack introduced restriction categories \cite{CockettLack},
and Cockett and Hofstra developed the theory of Turing categories
and related them to partial combinatory algebras through idempotent
splitting \cite[Theorem~4.9]{CockettHofstra}.
We use restriction categories, with binary
restriction products in the source recursion categories. These products
give the equivalent $P$-category structure
\cite[Proposition 4.5]{CockettLackIII}.

Cockett's later incompleteness theorem assumes a predicate expressing
program restriction order and produces a restriction idempotent on the
restriction terminal object that is neither zero nor the identity
\cite[Section~7.5, Theorem~7.14]{CockettComputability}.
Its conclusion concerns that restriction idempotent, whereas the results
here concern generated sentence theories and their consistent represented
extensions.

Di Paola and Heller proved recursive inseparability in dominical
recursion categories with a section for every morphism and binary
coproducts for which coproducts of total morphisms are total and the
near product distributes over coproducts
\cite[Theorem 8.15]{DPH}. Bernardi obtained categorical fixed-point and
inseparability results from precomplete equivalence relations, including PA-provable
equality of partial recursive functions among his examples
\cite[Section~3, pp.~125--128]{Bernardi}.
These results are recursion-theoretic antecedents of the present
incompleteness criteria.
Montagna also obtained effectively inseparable pairs in his syntactic
category $S$
\cite[Remark 3.1]{Montagna}.

Ramesh's introspective theories are essentially algebraic theories
extending the theory of finite-limit categories. Each model is equipped
with a homomorphism into an internal model of the same theory in its
underlying finite-limit category. From this interpretation he derives L\"ob's theorem and
guarded fixed-point results
\cite[Introduction, Definition~2.7, Theorem~4.19]{Ramesh2023}.

Here formula objects, initial algebras of finite derivations, and
conclusion images define internal deductive systems and their generated
sentence theories. The comparison and inseparability criteria apply to
the restriction idempotents representing these theories and their
refutations, yielding incompleteness throughout the stated consistent
represented extension classes. The criteria use categorical recursion
and inseparability; they do not postulate a provability operator
satisfying L\"ob's theorem or the Hilbert--Bernays derivability conditions.

Miller constructed a sequence that enumerates every real in $[0,1]$
computable uniformly from its representations \cite[Theorem~6.3]{Miller2004}.
We use the stronger index property stated immediately after that
theorem and recalled in Section~\ref{sec:miller}.
Bauer and Hanson use this construction in a parameterized realizability
topos with countable Dedekind reals \cite{BauerHanson}.
Section~\ref{sec:miller} uses their binary representations to construct a recursion category,
prove inseparability of $L_a,R_a$, and obtain essential incompleteness
for a theory of real closed fields.
